\documentclass[11pt,a4paper]{article}

% ── Packages ──────────────────────────────────────────────────────────────────
\usepackage[top=2.8cm, bottom=2.8cm, left=3cm, right=3cm]{geometry}
\usepackage[T1]{fontenc}
\usepackage[utf8]{inputenc}
\usepackage{lmodern}
\usepackage{microtype}
\usepackage[colorlinks=true, linkcolor=linkblue, urlcolor=accent,
            pdftitle={L1-09 Object Storage -- System Design Foundations},
            bookmarks=true]{hyperref}
\usepackage{xcolor}
\usepackage{titlesec}
\usepackage{enumitem}
\usepackage{booktabs}
\usepackage{tabularx}
\usepackage{tcolorbox}
\usepackage{fancyhdr}
\usepackage{amsmath}
\usepackage{parskip}
\usepackage{mdframed}
\usepackage{graphicx}
\usepackage{tikz}
\usepackage{setspace}
\usepackage{soul}
\usepackage{array}
\usepackage{longtable}

% ── Colors ────────────────────────────────────────────────────────────────────
\definecolor{primary}{HTML}{1A1A2E}
\definecolor{accent}{HTML}{C0392B}
\definecolor{highlight}{HTML}{1A5276}
\definecolor{linkblue}{HTML}{154360}
\definecolor{lightbg}{HTML}{F4F6F7}
\definecolor{warnbg}{HTML}{FEF9E7}
\definecolor{quizbg}{HTML}{EBF5FB}
\definecolor{termbg}{HTML}{EAF2FF}
\definecolor{curiobg}{HTML}{F5EEF8}
\definecolor{greenbg}{HTML}{EAFAF1}
\definecolor{notebg}{HTML}{FDF2E9}
\definecolor{accentlight}{HTML}{FADBD8}

% ── Section styling ───────────────────────────────────────────────────────────
\titleformat{\section}
  {\large\bfseries\color{highlight}}
  {\thesection.}{0.6em}{}
  [\vspace{2pt}\color{highlight!60}\rule{\linewidth}{0.6pt}]

\titleformat{\subsection}
  {\normalsize\bfseries\color{primary}}
  {\thesubsection}{0.5em}{}

\titleformat{\subsubsection}
  {\normalsize\itshape\color{primary!80}}
  {}{0em}{}

\titlespacing{\section}{0pt}{18pt}{8pt}
\titlespacing{\subsection}{0pt}{12pt}{4pt}
\titlespacing{\subsubsection}{0pt}{8pt}{2pt}

% ── Header/Footer ─────────────────────────────────────────────────────────────
\pagestyle{fancy}
\fancyhf{}
\fancyhead[L]{\small\color{highlight}\textbf{L1-09 --- Storage: Blob \& Object}}
\fancyhead[R]{\small\color{primary!70}System Design Interview Prep}
\fancyfoot[C]{\small\color{primary!60}\thepage}
\renewcommand{\headrulewidth}{0.4pt}
\renewcommand{\headrule}{\hbox to\headwidth{\color{highlight!40}\leaders\hrule height \headrulewidth\hfill}}

% ── Box environments ──────────────────────────────────────────────────────────
\tcbuselibrary{skins, breakable, hooks}

\newtcolorbox{termbox}[1]{
  enhanced, breakable,
  colback=termbg, colframe=highlight, coltitle=white,
  fonttitle=\bfseries\small, title={#1},
  left=8pt, right=8pt, top=6pt, bottom=6pt,
  attach boxed title to top left={yshift=-2.5mm, xshift=6mm},
  boxed title style={colback=highlight, sharp corners, left=4pt, right=4pt}
}

\newtcolorbox{industrybox}[1]{
  enhanced, breakable,
  colback=greenbg, colframe=highlight!50!black, coltitle=white,
  fonttitle=\bfseries\small, title={Industry Data: #1},
  left=8pt, right=8pt, top=6pt, bottom=6pt,
  attach boxed title to top left={yshift=-2.5mm, xshift=6mm},
  boxed title style={colback=highlight!60!black, sharp corners, left=4pt, right=4pt}
}

\newtcolorbox{trapbox}[1]{
  enhanced,
  colback=accentlight, colframe=accent, coltitle=white,
  fonttitle=\bfseries\small, title={Interview Trap: #1},
  left=8pt, right=8pt, top=6pt, bottom=6pt,
  attach boxed title to top left={yshift=-2.5mm, xshift=6mm},
  boxed title style={colback=accent, sharp corners, left=4pt, right=4pt}
}

\newtcolorbox{thinkbox}{
  enhanced,
  colback=notebg, colframe=accent!60,
  left=8pt, right=8pt, top=6pt, bottom=6pt,
  leftrule=3pt, rightrule=0pt, toprule=0pt, bottomrule=0pt,
}

\newtcolorbox{curiobox}[1]{
  enhanced, breakable,
  colback=curiobg, colframe=primary!60, coltitle=white,
  fonttitle=\bfseries\small, title={Q: #1},
  left=8pt, right=8pt, top=6pt, bottom=6pt,
  attach boxed title to top left={yshift=-2.5mm, xshift=6mm},
  boxed title style={colback=primary!70, sharp corners, left=4pt, right=4pt}
}

\newtcolorbox{quizbox}{
  enhanced, breakable,
  colback=quizbg, colframe=highlight, coltitle=white,
  fonttitle=\bfseries, title={Self-Quiz},
  left=10pt, right=10pt, top=8pt, bottom=8pt,
  attach boxed title to top left={yshift=-2.5mm, xshift=6mm},
  boxed title style={colback=highlight, sharp corners, left=4pt, right=4pt}
}

\newtcolorbox{answerbox}{
  enhanced, breakable,
  colback=lightbg, colframe=primary!40, coltitle=white,
  fonttitle=\bfseries\small, title={Model Answers},
  left=10pt, right=10pt, top=8pt, bottom=8pt,
  attach boxed title to top left={yshift=-2.5mm, xshift=6mm},
  boxed title style={colback=primary!60, sharp corners, left=4pt, right=4pt}
}

\newtcolorbox{insightbox}{
  enhanced,
  colback=primary!5, colframe=primary,
  left=8pt, right=8pt, top=6pt, bottom=6pt,
  leftrule=4pt, rightrule=0.4pt, toprule=0.4pt, bottomrule=0.4pt,
}

\setstretch{1.15}

% ── Document ──────────────────────────────────────────────────────────────────
\begin{document}

% ── Title page ────────────────────────────────────────────────────────────────
\begin{titlepage}
\vspace*{2cm}
\begin{center}
  {\color{primary}\rule{\linewidth}{3pt}}\\[12pt]
  {\huge\bfseries\color{highlight} L1-09}\\[6pt]
  {\Huge\bfseries\color{primary} Storage: Blob \& Object}\\[10pt]
  {\large\color{primary!70} The Infinite Shelf That Powers Video Streaming, File Sharing, and Every ML Pipeline on Earth}\\[20pt]
  {\color{primary}\rule{\linewidth}{1pt}}\\[16pt]
  {\normalsize\color{primary!70}
    \textit{Layer 1 --- Foundations} \quad\textbullet\quad
    \textit{System Design Interview Preparation} \quad\textbullet\quad
    \textit{2026 Edition}
  }\\[8pt]
  {\small\color{primary!60} Industry data sourced from AWS re:Invent, Netflix tech blog, Dropbox engineering, Cloudflare blog, Meta engineering, and AWS documentation (2022--2025)}
\end{center}
\vfill
\begin{insightbox}
\textbf{Why this lesson exists.} Every non-trivial system design interview eventually asks you to store something large: videos, user photos, model checkpoints, audit logs, backups. The instinctive answer---``just put it in the database''---fails immediately at scale. Understanding the three storage paradigms (block, file, object) and when each applies is what separates candidates who can design a Dropbox clone or a video upload pipeline from those who cannot. Object storage in particular is the invisible backbone of the modern internet: Netflix streams from S3, Instagram stores billions of photos in object storage, and every major ML training run reads its dataset from a bucket. This lesson gives you the mental model and the precise numbers to discuss this confidently in any interview.
\end{insightbox}
\vspace{1cm}
\tableofcontents
\end{titlepage}

% ─────────────────────────────────────────────────────────────────────────────
\section{Why This Exists: The Unstructured Data Problem}
% ─────────────────────────────────────────────────────────────────────────────

In the late 1990s, two things happened simultaneously in computing. First, disk drives got cheap enough that storing large files---images, audio, video---became economically feasible at consumer scale. Second, the web made it normal for millions of users to upload content. The combination created a problem that nobody had cleanly solved: how do you store billions of files, each potentially gigabytes in size, in a way that is cheap, durable, globally accessible, and operable without an army of storage engineers?

The existing answers were unsatisfying. Relational databases were designed for structured records, not raw bytes of video. File systems required a mountable volume, limiting scale to what a single machine's disk could hold. SAN (Storage Area Networks) and NAS (Network Attached Storage) scaled somewhat, but they were expensive, complex, and required specialized hardware. They were designed for the era of a few servers in a datacenter, not for a service that needed to store ten billion objects and serve them to users in São Paulo, Lagos, and Seoul simultaneously.

The engineering insight that changed everything was deceptively simple: if you give up the requirement that storage behave like a file system---no directory hierarchy, no mounts, no POSIX semantics---you can build something far more scalable. Instead of files-in-folders, you build a flat namespace of key-to-value mappings, each value being an arbitrary blob of bytes, each key being a string. You access objects via HTTP. You store the data redundantly across machines and datacenters. You make durability your primary guarantee. This is object storage.

Amazon launched S3 in March 2006. It was not the most featureful storage product ever built, but it was the most scalable. Every competitor since has either adopted the S3 API or built something compatible with it, because the mental model turned out to be exactly right: a flat, infinitely scalable, HTTP-accessible blob store.

\begin{industrybox}{S3 at Scale (AWS re:Invent 2023)}
As of 2023, Amazon S3 stores more than \textbf{350 trillion objects} and handles \textbf{100 million requests per second} at peak. The system has grown continuously since 2006 without ever experiencing a data loss event for correctly uploaded objects. The 11-nines durability figure (99.999999999\%) is not marketing; it is derived from the replication factor and failure-independence assumptions verified by AWS's own hardware failure rates. Source: AWS re:Invent 2023, keynote and storage deep-dive sessions.
\end{industrybox}

% ─────────────────────────────────────────────────────────────────────────────
\section{Three Storage Paradigms: Block, File, and Object}
% ─────────────────────────────────────────────────────────────────────────────

Before discussing object storage deeply, you must understand the full landscape. In every system design interview where storage comes up, you may be asked to justify your choice. Saying ``I'll use S3'' without being able to contrast it against block and file storage will cost you points.

\subsection{Block storage: raw bytes, maximum performance}

Block storage presents raw storage capacity to a computing host as a virtual disk. The host's operating system sees the storage as a block device---identical in concept to a physical hard drive or SSD---and formats it with whatever filesystem it wants (ext4, APFS, NTFS). The application on top of that host reads and writes as if it were talking to a local disk. The block storage system handles the low-level detail of where bytes are physically stored.

The defining characteristic of block storage is that \textbf{it has no concept of files or objects at the block storage layer}. It is a sequence of fixed-size blocks (typically 512 bytes or 4 KB) that can be read and written at any offset. This makes it ideal for workloads that need the lowest possible latency and full control over I/O patterns: databases, virtual machine images, and operating system volumes. A relational database like PostgreSQL uses block storage directly because it manages its own pages, write-ahead log, and buffer cache---it does not want an interposing filesystem making decisions about layout.

AWS Elastic Block Store (EBS) is the canonical cloud block storage service. An EBS volume attaches to a single EC2 instance (in most configurations) like a virtual drive. The AWS io2 Block Express tier guarantees up to 256,000 IOPS and 4,000 MB/s throughput with sub-millisecond latency. EBS volumes persist independently of the EC2 instance lifecycle---you can stop the instance, detach the volume, and reattach it to a different instance later. This is what makes it viable as database storage in the cloud.

\begin{thinkbox}
\textbf{Why can't block storage be attached to multiple instances simultaneously?} Because two independent processes writing to the same block device without coordination would corrupt each other's data. Block storage provides raw I/O without any locking or consistency protocol at the storage layer; that coordination must happen at a higher level. Clustered filesystems like OCFS2 or GFS2 exist precisely to solve this, but they add significant complexity and are only needed for specific HA database configurations. For most workloads, single-attachment block storage is exactly right.
\end{thinkbox}

\subsection{File storage: shared hierarchical namespace}

File storage adds a filesystem abstraction on top of shared storage. Multiple hosts can mount the same storage namespace simultaneously and read/write files using standard POSIX filesystem semantics: paths, directories, permissions, rename, and so on. This is what NFS (Network File System) and SMB/CIFS do, and it is what AWS EFS (Elastic File System) provides in the cloud.

The key advantage is multi-host access. If you have a fleet of application servers that all need to read and write a shared directory---legacy enterprise applications, media editing workstations, HPC compute nodes sharing input data---file storage is the natural fit. The shared namespace means every host sees the same view of the directory tree without any application-level coordination.

The tradeoff is scale and cost. File storage must maintain a coherent directory tree, metadata about every file (inodes, permissions, timestamps), and locking semantics when multiple writers touch the same data. This consistency overhead means file storage does not scale to billions of files the way object storage does. AWS EFS costs roughly \$0.30 per GB-month compared to S3 Standard at \$0.023 per GB-month---more than 10 times more expensive, because of the per-inode metadata and distributed locking machinery it must maintain.

File storage is the right answer when applications need POSIX filesystem semantics and cannot be modified: legacy applications that use \texttt{fopen()/fread()/fwrite()}, shared build artifact caches across CI nodes, or HPC jobs that read input data via mount points. It is the wrong answer for storing user-uploaded photos at Instagram scale.

\subsection{Object storage: infinite, cheap, HTTP-native}

Object storage is the storage model designed specifically for the internet era. It abandons the filesystem hierarchy entirely. There are no directories, no inodes, no mount points, and no locking. Instead, there are two constructs: \textbf{buckets} and \textbf{objects}.

A bucket is a globally named namespace within a storage region. An object is an arbitrary blob of bytes (from one byte to five terabytes in S3's case) identified by a string key within a bucket. You interact with objects via HTTP: a PUT request writes an object, a GET request retrieves it, a DELETE request removes it. That is the entire interface. There is no rename (you copy to a new key, delete the old one). There is no partial overwrite (you replace the entire object). These limitations are not bugs; they are what make object storage scale to trillions of objects.

\begin{termbox}{Object Storage Mental Model}
Think of object storage as an infinitely large dictionary: \texttt{bucket\_name + key $\rightarrow$ bytes + metadata}. The key is an arbitrary string. You can use slashes in keys (\texttt{photos/2024/january/img001.jpg}) to simulate a directory hierarchy, but the storage system treats this as a flat string. Object storage is not a filesystem with HTTP access; it is a different abstraction entirely that happens to be accessible via HTTP.
\end{termbox}

The three storage types side by side:

\begin{center}
\renewcommand{\arraystretch}{1.4}
\begin{tabular}{>{\bfseries}p{3.2cm}p{3.4cm}p{3.4cm}p{3.4cm}}
\toprule
 & \textbf{Block (EBS)} & \textbf{File (EFS)} & \textbf{Object (S3)} \\
\midrule
Access protocol & Block I/O (kernel) & NFS / SMB / POSIX & HTTP (REST API) \\
Hierarchy & Filesystem on top & Full directory tree & Flat key namespace \\
Concurrency & Single writer & Multi-host mount & Multi-client (eventual) \\
Latency & Sub-millisecond & Low ms & Tens of ms \\
Scalability & Per-volume (up to 64 TB) & Petabyte-scale & Essentially unlimited \\
Cost (AWS, \$/GB-mo) & \$0.08--0.125 (gp3) & \$0.30 (standard) & \$0.023 (standard) \\
Best for & Databases, OS volumes & Shared apps, HPC & Media, backups, ML data \\
\bottomrule
\end{tabular}
\end{center}

% ─────────────────────────────────────────────────────────────────────────────
\section{S3 Architecture: How Amazon Built an Infinite Shelf}
% ─────────────────────────────────────────────────────────────────────────────

Understanding S3's internal architecture is not just academic---it directly informs how you use it correctly, why certain limitations exist, and how to design systems around it.

\subsection{Buckets: namespace and policy boundary}

A bucket is the top-level organizational unit in S3. Every bucket has a globally unique name (across all AWS accounts and all regions), is created in a specific AWS region, and serves as the unit for access control policies, versioning configuration, replication rules, and lifecycle policies. You cannot transfer a bucket between regions; if you need data in a different region, you use Cross-Region Replication.

The global uniqueness requirement has a practical implication: if you try to create a bucket named \texttt{my-app-data} and that name is already taken by anyone in the world, you cannot use it. This occasionally surprises teams that try to programmatically create buckets with predictable names and find them already claimed. The convention is to include your AWS account ID or a random suffix in bucket names to ensure uniqueness.

Buckets do not limit the number of objects they contain---a single bucket can hold trillions of objects with no performance degradation if designed correctly. They do not limit total size. The limits that matter are per-object (5 TB) and request rate (which we will address in the presigned URL and multipart sections).

\subsection{Objects: key, value, metadata, and ETags}

Every object in S3 consists of four components. The \textbf{key} is the object's name within the bucket, a string up to 1,024 bytes. The \textbf{value} is the data itself, from 0 bytes to 5 TB. The \textbf{metadata} is a collection of key-value string pairs attached to the object: system metadata (Content-Type, Content-Length, ETag, Last-Modified) and up to 2 KB of user-defined metadata. The \textbf{ETag} is an MD5 checksum (or a compound hash for multipart uploads) that clients can use for cache validation and integrity checking.

The slash-as-delimiter convention is worth understanding precisely. If you store objects with keys like \texttt{videos/2024/intro.mp4} and \texttt{videos/2024/chapter1.mp4}, the S3 console and list APIs will display \texttt{videos/} as a ``folder.'' But S3 does not store a folder object; the slash is just part of the key string. The ListObjectsV2 API lets you specify a \texttt{Delimiter} parameter so that responses group keys with a common prefix---this is how the illusion of folders is implemented. Knowing this prevents a subtle trap: never design application logic that depends on folder operations having any special semantics in S3, because they do not.

\begin{thinkbox}
\textbf{Why does the key design matter for performance?} S3 partitions objects across its internal storage nodes based on the key prefix. If you name all your objects with a date prefix (e.g., \texttt{2024-01-15-image-001.jpg}), all objects created on that day hash to the same partition, and you hit request rate limits for that partition faster. AWS's guidance since 2018: use randomized prefixes (a hash of the object ID, not the date) to distribute load across partitions. This matters when your upload pipeline generates tens of thousands of objects per second.
\end{thinkbox}

\subsection{Versioning: protecting against accidental deletion}

S3 versioning is a bucket-level setting. When enabled, every PUT, COPY, or DELETE request creates a new version of the object rather than overwriting it. Each version gets a unique version ID string. A DELETE request without a version ID inserts a delete marker---the object appears gone to regular GET requests, but all previous versions remain. To truly delete an object, you must delete all its versions and the delete marker.

Versioning is invaluable for two scenarios. First, accidental deletion: if a bug or errant script deletes objects, you recover by removing the delete markers. Second, concurrent writes from different clients: if two processes PUT to the same key simultaneously, versioning ensures neither write is lost---you get two versions and can reconcile them in your application.

The cost of versioning is storage: every version is stored and billed independently. A file that changes daily accumulates 365 versions per year, each consuming storage. S3 lifecycle policies let you configure automatic deletion of old versions (e.g., ``keep the last 5 versions, delete the rest after 90 days'') to manage costs.

\subsection{Storage classes: paying for what you actually need}

S3 is not a single storage tier---it is a family of storage classes with different cost and retrieval latency tradeoffs. The key insight is that the vast majority of data in most systems is \textbf{rarely accessed after a brief initial window}. A video uploaded today is watched heavily in the first week, then accessed infrequently for months, then almost never. Storing it in the same tier forever wastes money.

\begin{center}
\renewcommand{\arraystretch}{1.4}
\begin{tabular}{>{\bfseries}p{3.8cm}p{1.8cm}p{2.5cm}p{4.2cm}}
\toprule
\textbf{Storage Class} & \textbf{Cost (approx.)} & \textbf{Retrieval time} & \textbf{Use case} \\
\midrule
S3 Standard & \$0.023/GB-mo & Milliseconds & Active content, hot data \\
S3 Standard-IA & \$0.0125/GB-mo & Milliseconds & Infrequent access, DR copies \\
S3 One Zone-IA & \$0.01/GB-mo & Milliseconds & Reproducible data, lower durability \\
S3 Intelligent-Tiering & \$0.023--0.004/GB & Milliseconds & Unknown or changing access patterns \\
S3 Glacier Instant & \$0.004/GB-mo & Milliseconds & Archives needing rapid access \\
S3 Glacier Flexible & \$0.0036/GB-mo & Minutes--hours & Compliance archives, infrequent retrieval \\
S3 Glacier Deep Archive & \$0.00099/GB-mo & 12--48 hours & Long-term cold archive, lowest cost \\
\bottomrule
\end{tabular}
\end{center}

S3 Lifecycle policies automate the movement of objects between storage classes. A typical media platform rule: move objects to Standard-IA after 30 days, Glacier Instant Retrieval after 180 days, and Deep Archive after 2 years. This can reduce storage costs by 80--90\% for data that is almost never accessed after the initial upload window.

\begin{industrybox}{Netflix and Glacier for Cold Archive (Netflix Tech Blog, 2023)}
Netflix uses S3 Glacier for master copies of content that is not currently licensed for streaming but must be retained for future re-licensing negotiations. As of 2023, Netflix stores petabytes of raw production footage in Glacier Deep Archive at roughly \$0.001/GB-month, compared to storing it on-premise on LTO tape. The transition to cloud archive storage reduced operational complexity (no tape robot maintenance, no media degradation) while cutting per-GB cost below what physical tape infrastructure achieved at their scale.
\end{industrybox}

% ─────────────────────────────────────────────────────────────────────────────
\section{Presigned URLs: Direct Client Uploads Without Your Backend}
% ─────────────────────────────────────────────────────────────────────────────

One of the most important patterns in object storage, and one of the most frequently misunderstood in interviews, is the presigned URL. Understanding why it exists is more important than knowing the API.

\subsection{The problem: large file uploads through your backend}

Consider the naive architecture for a user uploading a video: the client sends the file to your API server, which buffers it, then uploads it to S3. This design has three serious problems at scale.

First, \textbf{bandwidth consumption on your servers}. A 1 GB video upload consumes 1 GB of inbound bandwidth on your API server and 1 GB of outbound bandwidth from your server to S3. You are paying twice for the same transfer: once for client-to-server and once for server-to-S3. Your server fleet needs to be sized to handle peak upload bandwidth, which for a consumer video platform can be enormous.

Second, \textbf{server memory and CPU}. Buffering large files in memory or on disk while uploading them to S3 adds memory pressure and disk I/O to servers that should be focused on business logic. A 100-MB video upload that takes 60 seconds ties up a server thread or worker process for a full minute.

Third, \textbf{single point of failure}. If your API server goes down mid-upload, the client must restart the entire transfer from scratch. There is no mechanism in this architecture to resume.

\subsection{The solution: the client uploads directly to S3}

Presigned URLs solve this by letting your backend sign a time-limited URL that grants the bearer the right to perform one specific S3 operation on one specific object---typically a PUT to upload a file. The client receives this URL and uses it to upload directly to S3 without the data ever touching your server.

\begin{termbox}{Presigned URL}
A presigned URL is an S3 URL that includes a cryptographic signature in its query parameters. The signature is created by your backend using your AWS credentials and encodes: the bucket, the object key, the HTTP method (GET or PUT), and an expiration time. Any HTTP client that has this URL can perform the signed operation within the expiration window, without needing AWS credentials of their own. After the URL expires, the signature is invalid and S3 rejects the request.
\end{termbox}

The flow for a presigned upload:

\begin{enumerate}[leftmargin=1.5em, itemsep=4pt]
  \item Client requests an upload URL from your API: \texttt{POST /api/upload-url \{filename: "video.mp4", size: 104857600\}}
  \item Your backend validates the request (is this user allowed to upload? is the file type acceptable?), generates a presigned PUT URL pointing to S3, and returns it to the client.
  \item Client uses the presigned URL to PUT the file directly to S3. No data flows through your server.
  \item S3 stores the file and optionally triggers an S3 event notification (to SQS, SNS, or Lambda) informing your system that the upload is complete.
  \item Your backend receives the S3 event and updates its database to mark the upload as finalized.
\end{enumerate}

The bandwidth saving is real and substantial. Dropbox's engineering team published in 2014 that switching to direct-to-S3 uploads eliminated an entire tier of upload proxy servers from their architecture. At their scale (hundreds of millions of uploads per day), eliminating the proxy hop reduced upload latency by 20--30\% and cut infrastructure costs significantly.

\subsection{Presigned URL TTL and security}

Presigned URLs must have an expiration time. AWS allows up to 7 days (604,800 seconds) for URLs signed with long-term credentials, and up to the expiration of the underlying IAM session for URLs signed with temporary credentials (STS-assumed roles, which are typically scoped to 1--12 hours).

For upload URLs, best practice is a short TTL---typically 5 to 15 minutes. The URL should be requested immediately before the upload begins, not cached. If the URL leaks (e.g., logged in a client-side error report), its short TTL limits the damage window. For download URLs serving sensitive content, TTL should match your business policy: a URL for a healthcare document might expire in 1 minute; a URL for a shared photo might last an hour.

\begin{trapbox}{Treating Presigned URLs as Permanent Links}
A common interview mistake is to say ``generate a presigned URL when the file is uploaded and store it in your database.'' This is wrong for two reasons. First, presigned URLs expire---within hours or days, your stored URL becomes a 403. Second, storing signed URLs in your database means your database becomes the source of truth for authorization, but S3 is the one actually enforcing the expiration. The correct model: store only the S3 bucket name and object key in your database. Generate presigned URLs on demand, at request time, using your IAM credentials. This way authorization is always fresh and your signing keys can be rotated without invalidating stored data.
\end{trapbox}

% ─────────────────────────────────────────────────────────────────────────────
\section{Multipart Upload: Handling Files Larger Than Memory}
% ─────────────────────────────────────────────────────────────────────────────

S3 imposes a 5 GB limit on single-request PUT uploads. But more importantly, even well below that limit, large file uploads over unreliable networks (mobile connections, intercontinental paths) fail frequently. When an 800 MB upload fails after transferring 750 MB, restarting from zero is both wasteful and user-hostile. Multipart upload exists to solve this.

\subsection{How multipart upload works}

Multipart upload is a three-phase protocol. In phase one, the client initiates the multipart upload, receiving a unique Upload ID from S3. In phase two, the client splits the file into parts (each between 5 MB and 5 GB, except the final part which can be smaller) and uploads each part independently as a numbered PUT request. S3 stores each part separately and returns an ETag for each. Parts can be uploaded concurrently---the client can open multiple parallel streams, saturating the available bandwidth. In phase three, the client sends a completion request listing all part numbers and their ETags in order. S3 assembles the final object from the stored parts.

\begin{termbox}{Multipart Upload Constraints}
Minimum part size: \textbf{5 MB} (except the last part). Maximum part size: \textbf{5 GB}. Maximum number of parts: \textbf{10,000}. This means the maximum object size via multipart upload is $10{,}000 \times 5\text{ GB} = 5\text{ TB}$, which is also the maximum S3 object size. AWS recommends using multipart upload for any object larger than \textbf{100 MB}. For objects smaller than 5 MB, single-part PUT is the only option (5 MB is the minimum part size, so you cannot initiate a multipart upload for a 3 MB file).
\end{termbox}

The key operational insight is that multipart upload is resumable. If a network failure interrupts part 7 of 20, the client can retry part 7 independently---parts 1 through 6 are already stored at S3 and do not need to be re-sent. The partial parts remain stored in S3 until either the upload is completed or aborted. Crucially, \textbf{you are billed for partial multipart uploads that are never completed}. This is a real operational cost trap: if a client initiates a multipart upload, uploads some parts, and then abandons it (due to a crash or navigation away from the page), those parts accrue storage charges indefinitely unless cleaned up. S3 lifecycle policies have an \texttt{AbortIncompleteMultipartUpload} rule specifically for this scenario---set it to clean up incomplete uploads after a sensible number of days (7 is typical for most consumer-facing applications).

\begin{industrybox}{Parallel Multipart Upload Throughput (AWS Documentation, 2024)}
A client on a 10 Gbps network uploading a 10 GB object as 100 parts of 100 MB each, with 10 concurrent part uploads, can achieve near-line-rate throughput to S3. AWS internal benchmarks show that parallelizing part uploads is the single most effective technique for maximizing S3 upload throughput from a single host. The AWS SDK's built-in Transfer Manager automatically uses multipart upload for files above a configurable threshold (default: 8 MB) and manages parallelism transparently.
\end{industrybox}

% ─────────────────────────────────────────────────────────────────────────────
\section{S3 Consistency Model: What Changed in 2020}
% ─────────────────────────────────────────────────────────────────────────────

Before December 2020, S3 offered \textbf{read-after-write consistency} for new objects only. If you PUT an object for the first time, you could immediately GET it and receive the new data. But for overwrite PUTs and DELETEs, S3 offered only \textbf{eventual consistency}: after overwriting an object, a subsequent GET might still return the old version for some brief window. This was a well-known caveat that complicated S3-based architectures for years.

In December 2020, AWS announced that S3 now provides \textbf{strong read-after-write consistency for all operations}---new PUTs, overwrite PUTs, and DELETEs---at no additional cost. A GET immediately following any successful PUT or DELETE is guaranteed to reflect the latest state.

\begin{thinkbox}
\textbf{How did AWS achieve this without sacrificing performance?} The change required restructuring S3's internal metadata storage layer. AWS added a ``witness'' step in the write path: a write is not acknowledged as successful until it is durably recorded in a way that all subsequent reads can observe. The engineering challenge was making this fast enough that it did not degrade write throughput. AWS's approach---detailed in their 2023 OSDI paper on S3 internals---involved using a variant of Paxos for metadata coordination combined with read-repair on the hot path.
\end{thinkbox}

This matters for interview answers. If you are describing an architecture built before 2020, or if you encounter documentation referencing eventual consistency for overwrites, it is now out of date. The current guarantee is strong consistency for all standard GET, PUT, and DELETE operations. However, one important caveat remains: S3 is not a database. Strong consistency for individual object operations does not mean serializable transactions across multiple objects. If you need atomic multi-object updates, S3 alone cannot provide that; you need a coordinating database or transaction log.

% ─────────────────────────────────────────────────────────────────────────────
\section{Durability, Replication, and the Eleven Nines}
% ─────────────────────────────────────────────────────────────────────────────

S3's headline durability figure is \textbf{99.999999999\%} (eleven nines) for S3 Standard objects. This number sounds like marketing, but it has a concrete meaning: for any given object, the expected probability of losing it in a given year is $10^{-11}$. If you store 10 million objects for 10 years, you can statistically expect to lose approximately one of them due to hardware failure over that entire period.

This durability is achieved through two mechanisms. First, S3 stores every object across \textbf{at least three Availability Zones} (physically separate datacenters within a region, with independent power, network, and cooling). A single AZ failure---a power outage, a network partition, a flooding event---does not cause data loss because the other two AZs retain the data. Second, within each AZ, S3 stores objects on multiple independent hardware nodes and uses erasure coding for larger objects. Even if individual disks fail (which happens constantly in a fleet of hundreds of thousands of drives), the erasure coding can reconstruct the lost data.

S3 One Zone-IA breaks this model deliberately: it stores objects in only one AZ to reduce cost by 20\%. The tradeoff is that an AZ failure could destroy those objects. One Zone-IA is appropriate for data that can be regenerated cheaply (thumbnail images derived from originals, processed derivatives of raw data). It is never appropriate for primary copies of user-uploaded content.

\begin{termbox}{Cross-Region Replication (CRR)}
CRR is an S3 feature that automatically replicates objects from a source bucket in one AWS region to a destination bucket in another region. It is asynchronous---the replication typically completes within seconds to minutes but has no strict SLA. CRR is used for: (1) disaster recovery, so that a full-region AWS outage does not cause data unavailability; (2) latency reduction, keeping copies geographically close to users in different regions; (3) data sovereignty compliance, ensuring a copy exists in a specific jurisdiction. Replication does not change the source object---it only creates additional copies.
\end{termbox}

The availability SLA for S3 Standard is 99.99\% per month (about 52 minutes of downtime per year). This is distinct from durability: availability measures whether requests succeed, while durability measures whether data exists. S3 can experience brief periods where API requests fail (an availability event) while the underlying data remains perfectly intact and durable.

% ─────────────────────────────────────────────────────────────────────────────
\section{CDN Integration: Getting Bytes to Users Fast}
% ─────────────────────────────────────────────────────────────────────────────

Object storage durability and scale are necessary but not sufficient for a great user experience. An S3 bucket in \texttt{us-east-1} can serve objects reliably, but a user in Singapore retrieving a 500 MB video from a Virginia datacenter will experience high latency and potentially poor throughput due to the round-trip time across the Pacific Ocean. The solution is a Content Delivery Network (CDN).

A CDN places edge caches at Points of Presence (PoPs) distributed globally---hundreds of locations. When a user requests an object, the CDN routes the request to the nearest PoP. If the PoP has the object cached (a cache hit), it serves it locally with sub-10ms latency. If not (a cache miss), the PoP fetches the object from the S3 origin, caches it, and serves the user. Subsequent users at the same PoP get the cached version.

The canonical integration is Amazon CloudFront with S3 as the origin. The architecture has three elements: the S3 bucket stores the canonical, durable copy of every object; CloudFront serves as the public-facing entry point; and an Origin Access Control (OAC) policy ensures that S3 rejects direct requests that do not come through CloudFront, preventing cache bypassing.

\begin{industrybox}{Netflix Open Connect CDN (Netflix Tech Blog, 2022)}
Netflix operates its own CDN, Netflix Open Connect, which embeds purpose-built appliances in ISP networks. As of 2022, Open Connect serves more than \textbf{15 petabytes of data per day} from approximately \textbf{17,000 servers} at thousands of ISP locations worldwide. The master copies of Netflix content---original video files, transcoded renditions---live in S3. Open Connect appliances are pre-populated nightly with anticipated content based on regional popularity predictions. A user in Tokyo watching a popular US show may receive every byte from an appliance inside their ISP's own building. Source: Netflix Tech Blog, ``Open Connect Continues Scaling'' (2022).
\end{industrybox}

For most systems without Netflix's scale to operate their own CDN, the S3 plus CloudFront (or S3 plus Cloudflare) combination achieves most of the benefit. S3 Transfer Acceleration is an alternative for upload scenarios: it routes uploads through the CloudFront edge network rather than directly to S3, improving upload throughput from distant clients by 50--500\% in AWS's measurements.

% ─────────────────────────────────────────────────────────────────────────────
\section{Decision Framework: Block vs File vs Object}
% ─────────────────────────────────────────────────────────────────────────────

In interviews, you will be asked to justify storage choices. The decision matrix below captures the key signals.

Use \textbf{block storage} when: (a) a relational database is involved---PostgreSQL, MySQL, Oracle all require block storage for their data files and WAL; (b) the workload is latency-sensitive and requires sub-millisecond I/O---trading systems, OLTP at high concurrency; (c) the operating system's boot volume. Block storage is the right answer whenever you are running software that manages its own I/O and cannot be modified to use an HTTP API.

Use \textbf{file storage} when: (a) multiple servers must simultaneously access a shared directory with POSIX semantics; (b) the application is a legacy system that uses \texttt{fopen()} and cannot be modified; (c) shared configuration files, shared build caches, or home directories in a multi-server environment. File storage is typically a compatibility choice---you reach for it when you cannot use object storage because the software requires a filesystem interface.

Use \textbf{object storage} when: (a) you are storing user-generated content: photos, videos, documents, audio; (b) the data is immutable after initial write (or infrequently updated); (c) you need global access via HTTPS without VPN or private network; (d) cost is a primary concern and you need to store more than a few terabytes; (e) you are building a data pipeline that ingests, processes, and archives large files (ML training data, raw event logs, batch job outputs). Object storage is the default for the modern internet.

\begin{trapbox}{Using a Database as a File Store}
A frequent interview mistake: ``We'll store user uploads in the database as BLOBs.'' For small files (profile pictures under a few KB), this is arguably acceptable. For anything larger, it is wrong. Databases are optimized for transactional record storage and OLTP query patterns. Storing megabytes-to-gigabytes in database rows bloats table sizes, increases backup times, makes vacuuming and compaction slower, and prevents the database from being optimized for its actual workload. The correct answer: store the file in S3, store the S3 key in the database. The database record is a few bytes; the file is in purpose-built storage.
\end{trapbox}

% ─────────────────────────────────────────────────────────────────────────────
\section{Critical Numbers Reference}
% ─────────────────────────────────────────────────────────────────────────────

\begin{center}
\renewcommand{\arraystretch}{1.5}
\begin{tabular}{p{5.5cm}p{4cm}p{4cm}}
\toprule
\textbf{Metric} & \textbf{Value} & \textbf{Source} \\
\midrule
S3 max object size (single PUT) & 5 GB & AWS Documentation \\
S3 max object size (multipart) & 5 TB & AWS Documentation \\
S3 max number of parts & 10,000 & AWS Documentation \\
S3 min part size (multipart) & 5 MB (last part: any) & AWS Documentation \\
S3 max key length & 1,024 bytes & AWS Documentation \\
S3 Standard durability & 99.999999999\% (11 nines) & AWS SLA \\
S3 Standard availability & 99.99\% per month & AWS SLA \\
S3 Standard-IA availability & 99.9\% per month & AWS SLA \\
S3 Standard cost & \$0.023/GB-month & AWS Pricing (2025) \\
S3 Glacier Deep Archive cost & \$0.00099/GB-month & AWS Pricing (2025) \\
S3 request rate & Up to 3,500 PUT/s per prefix & AWS Documentation \\
S3 request rate & Up to 5,500 GET/s per prefix & AWS Documentation \\
S3 total objects stored (2023) & 350 trillion & AWS re:Invent 2023 \\
S3 peak requests (2023) & 100 million/second & AWS re:Invent 2023 \\
Presigned URL max TTL (permanent) & 7 days & AWS Documentation \\
Multipart recommended threshold & 100 MB & AWS Best Practices \\
\bottomrule
\end{tabular}
\end{center}

% ─────────────────────────────────────────────────────────────────────────────
\section{System Design Application}
% ─────────────────────────────────────────────────────────────────────────────

\subsection{Application 1: Dropbox-style file storage}

The canonical object storage interview question is a cloud file storage service like Dropbox. The storage architecture unfolds in layers. User files are stored as objects in S3, with the S3 key derived from the user's account ID and the file's content hash (enabling deduplication: two users who store an identical file share one S3 object). A metadata database (PostgreSQL or DynamoDB) stores the mapping from user namespace (folder structure, file names) to S3 keys, along with versioning history. Uploads use the presigned URL pattern: the client requests an upload URL from the API, uploads directly to S3, and the API updates the metadata database on S3 event completion.

The deduplication insight is worth elaborating. Dropbox published in 2014 that block-level deduplication reduces their stored data by approximately 30\%. The trick: before uploading, the client computes a hash of the file (or of each block for large files). The client asks the server ``do you already have hash X?'' If yes, the server simply records a pointer to the existing S3 object. No upload necessary. This is the core of Dropbox's Smart Sync architecture.

\begin{industrybox}{Dropbox Storage Architecture (Dropbox Engineering Blog, 2016)}
Dropbox migrated from AWS S3 to a custom-built storage system called Magic Pocket in 2016, storing the majority of their data in-house to reduce costs. However, the architecture they describe is instructive: data is stored as fixed-size 4 MB blocks, each block content-addressed by its SHA-256 hash, with deduplication at the block level. The metadata layer maps user files to lists of block hashes. As of 2016 the system stored more than \textbf{500 petabytes} of user data. Most companies without Dropbox's scale would use S3 with the same logical architecture.
\end{industrybox}

\subsection{Application 2: Video upload and processing pipeline}

A video platform like YouTube or TikTok needs to accept raw uploads (often large and from mobile clients), transcode them into multiple resolutions and formats, and serve them globally with low latency. The object storage layer is central to all three stages.

For upload, the client uses a presigned URL or, for very large files (4K raw footage can be tens of gigabytes), a presigned multipart upload. S3 stores the raw video as a single object (or as reassembled multipart parts). An S3 event notification fires to an SQS queue when the upload completes. A transcoding fleet reads the raw video from S3, produces multiple output files (240p, 480p, 1080p, an HLS manifest, thumbnail images), and writes each output to S3. The final outputs are served via CloudFront.

The key design points are: the raw video and transcoded outputs live in separate S3 buckets (the raw bucket is private; the output bucket is served via CDN). The transcoding pipeline is triggered by event notifications, making it event-driven and decoupled from the upload path. Lifecycle policies on the raw bucket move older raw files to Glacier after 90 days, since the transcoded versions are what users actually consume.

\subsection{Application 3: ML training data store}

Machine learning pipelines read enormous amounts of data: training datasets for large models can be hundreds of terabytes or petabytes of images, text, or video. Object storage is the de facto home for ML training data, and the architecture has specific patterns.

Training data is stored in S3 as individual files (one JPEG per image, one CSV per shard) or in container formats (TFRecord, Parquet, WebDataset) that pack thousands of records per file. The container format is critical for performance: reading ten thousand individual 50KB images requires ten thousand separate S3 GET requests (each with tens of milliseconds of latency) while reading ten large Parquet files requires ten requests. The AWS documentation for SageMaker recommends Pipe Mode or FSx for Lustre as a caching layer between S3 and training instances for latency-sensitive training jobs.

\begin{industrybox}{Meta's ML Training Data Storage (Meta Engineering, 2023)}
Meta's training infrastructure stores petabytes of training data for its recommendation models and LLMs in a distributed storage system they call Manifold, built on top of blob storage. Training jobs read data at up to \textbf{1 TB/s aggregate} across thousands of GPUs. To achieve this throughput, they use a multi-level caching architecture: SSDs on the training nodes, a distributed flash cache tier, and S3-compatible blob storage as the persistent backing store. Direct reads from S3 alone cannot saturate modern GPU training clusters---the caching layer is essential.
\end{industrybox}

% ─────────────────────────────────────────────────────────────────────────────
\section{Curiosity Corner}
% ─────────────────────────────────────────────────────────────────────────────

\begin{curiobox}{Why does S3 have a globally unique bucket namespace? Couldn't AWS just scope buckets to accounts?}
The global bucket namespace is a design decision from 2006 that AWS has never been able to change without breaking existing deployments. The reason it was chosen originally: S3 bucket names are used as subdomain prefixes in the virtual-hosted-style URL format (\texttt{mybucket.s3.amazonaws.com}). For DNS to work cleanly, bucket names must be globally unique, just as domain names are globally unique. AWS could have used account-scoped names with path-style URLs (\texttt{s3.amazonaws.com/accountid/bucketname}) but chose subdomain-style for cleanliness and CDN compatibility. The cost is namespace squatting---popular names are grabbed early by anyone with an AWS account, and organizations regularly find their preferred bucket names already taken. AWS has maintained backward compatibility at the expense of this ergonomic annoyance for nearly 20 years.
\end{curiobox}

\begin{curiobox}{What is Cloudflare R2 and why does it matter competitively?}
Cloudflare R2, launched in general availability in September 2022, is an S3-compatible object storage service with one critical differentiator: \textbf{zero egress fees}. AWS S3 charges \$0.09/GB for data transferred out to the internet. For companies streaming large amounts of data---video platforms, software distribution networks, data export services---egress fees often exceed storage fees. R2's storage pricing is comparable to S3 Standard at \$0.015/GB-month, but the \$0 egress cost makes it dramatically cheaper for high-bandwidth workloads. As of 2024, Cloudflare reported that R2 stores more than \textbf{10 billion objects} for its customers. The existence of R2 forced a conversation in the industry about egress fee pricing, and in 2024 AWS, Google, and Microsoft jointly announced they would waive egress fees for customers switching cloud providers---a direct response to R2's competitive pressure.
\end{curiobox}

\begin{curiobox}{How does S3 achieve 11-nines durability? Is it just replication?}
Simple 3x replication across AZs would give you approximately 9 nines, not 11. To reach 11 nines, S3 uses a combination of three techniques. First, cross-AZ replication: every object is stored in at least three geographically separated AZs. Second, erasure coding within each AZ: larger objects are broken into fragments with additional parity fragments (similar to RAID-6) so that even if several storage nodes within an AZ fail simultaneously, data can be reconstructed. Third, continuous background integrity checking: S3 runs a continuous process that reads stored objects, verifies checksums, and repairs bit-rot or silent corruption. This is what the 11 nines actually requires---not just redundancy, but active monitoring and repair, because hardware at AWS's scale experiences constant low-level failures (bit flips, disk sectors going bad, firmware bugs). The durability figure is the result of modeling these failure rates against the replication and repair mechanisms.
\end{curiobox}

\begin{curiobox}{Can I use S3 as a database? What are its query limitations?}
S3 Select and Athena exist precisely because people want SQL-like queries over data in S3. S3 Select is a server-side filtering API: given a CSV or JSON or Parquet object in S3, you can push a SQL WHERE clause to S3 and receive only the matching rows rather than the full object. This reduces egress and client-side processing. Amazon Athena is a serverless query engine that lets you run full SQL queries over S3 data using Presto under the hood, treating S3 prefixes as table partitions. Athena is used extensively for log analysis, business intelligence on raw event data, and ad hoc exploration of S3-resident datasets. The critical limitation is that neither S3 Select nor Athena provides OLTP semantics---there are no indexes in the traditional sense, no row-level locks, and queries scan data rather than performing key lookups. Query latency is measured in seconds to minutes, not milliseconds. For transactional workloads, you need a real database. For analytical workloads on large S3 datasets, Athena is cost-effective and powerful.
\end{curiobox}

\begin{curiobox}{How does Instagram store billions of photos? Is it all S3?}
Instagram's photo storage architecture is well-documented from their engineering blog posts. Early Instagram (2010--2012) stored photos on S3 with metadata in PostgreSQL. As they scaled to hundreds of millions of photos per day, they built a caching layer (using Memcache) in front of S3 for hot photos. In 2014, they open-sourced their approach to serving images efficiently: the key insight was that the majority of views go to a small fraction of photos (the viral ones), so an aggressive in-memory cache can absorb most traffic before it hits S3. By 2019, Instagram was storing more than \textbf{100 billion photos}. They maintain the original uploaded photo and multiple derived thumbnails, each stored as separate S3 objects. The metadata (which user uploaded it, who liked it, caption text) is in Cassandra and PostgreSQL. The photo bytes themselves are in S3 (or Facebook's internal equivalent, called Haystack for small photos and f4 for larger media). The S3 API compatibility layer means the application code is portable even if the underlying storage implementation changes.
\end{curiobox}

\begin{curiobox}{What happens when you delete objects from S3 versioning-enabled buckets?}
This is a subtlety that trips up many engineers in production. When you issue a DELETE on a versioned S3 bucket without specifying a version ID, S3 does not delete the object---it inserts a \textbf{delete marker} as the latest version. A GET request then returns 404, as if the object is gone, but all previous versions still exist and are still billed. To truly delete the object, you must delete each version individually, including the delete marker. If you accidentally ``delete'' a versioned object and need to recover it, simply delete the delete marker---the previous version becomes the current one. This behavior is the reason why enabling versioning can lead to surprisingly high storage bills: deleted files accumulate delete markers, and old versions accumulate silently beneath them. The lifecycle rule \texttt{NoncurrentVersionExpiration} automatically deletes versions that have been superseded, and \texttt{ExpiredObjectDeleteMarker} cleans up orphaned delete markers after all versions beneath them have been removed.
\end{curiobox}

% ─────────────────────────────────────────────────────────────────────────────
\section{Self-Quiz}
% ─────────────────────────────────────────────────────────────────────────────

\newpage

\begin{quizbox}

\textbf{Question 1.} You are designing a video platform where users upload raw 4K footage that may be 20--50 GB in size from mobile devices on variable connections. Walk me through the upload architecture. What storage mechanism do you use, and why?

\medskip
\textbf{Question 2.} A junior engineer on your team proposes storing user profile photos in the users table in PostgreSQL as binary columns. The system expects 50 million users and profile updates happening 2 million times per day. What is wrong with this proposal, and what do you recommend instead?

\medskip
\textbf{Question 3.} Your S3-based media platform is seeing users in Asia experience 3--5 second load times for images stored in \texttt{us-east-1}. The images change infrequently after initial upload. What do you change, and what are the tradeoffs?

\medskip
\textbf{Question 4.} You are storing compliance audit logs that must be retained for 7 years and retrieved at most once per year during audits. The logs are currently costing \$50,000/month on S3 Standard. How do you reduce that cost without violating the retention requirement, and what retrieval implications must your team accept?

\medskip
\textbf{Question 5.} Describe the S3 consistency model as it stands today. Then tell me: if two clients simultaneously PUT to the same S3 key, what happens? Is that a problem for your design, and how would you address it if it is?

\end{quizbox}

\newpage

\begin{answerbox}

\textbf{Answer 1.} Use S3 with multipart upload via presigned URLs. The client requests a presigned multipart upload initiation from your API, which returns an Upload ID. The client splits the file into parts (100--500 MB each, e.g., 100 parts for a 20 GB file) and uploads each part directly to S3 in parallel---no data flows through your API servers, eliminating bandwidth costs on your backend. Parts can be uploaded concurrently saturating the client's available bandwidth. If a mobile connection drops during part 7 of 100, only part 7 needs to be retried---parts 1--6 are already stored at S3. When all parts complete, the client sends a completion request. Your API receives an S3 event notification and updates the database with the final object key. Set lifecycle rules to abort incomplete multipart uploads after 7 days to avoid billing for abandoned uploads.

\medskip
\textbf{Answer 2.} Storing binary files in PostgreSQL columns bloats row size, inflates table scan times, stresses the WAL and autovacuum process, makes backups slower (the entire database dump includes every photo), prevents PostgreSQL from being optimized for transactional operations, and makes CDN delivery impossible without proxying every image through your API servers. The correct design: store photos as S3 objects using a key based on the user ID and a content hash. Store only the S3 key string (a few bytes) in the users table. On retrieval, generate a presigned URL (for private photos) or serve via CloudFront (for public photos). This separates transactional data from blob data cleanly.

\medskip
\textbf{Answer 3.} Place CloudFront (or another CDN) in front of the S3 bucket. CloudFront has Points of Presence throughout Asia---users in Tokyo, Singapore, and Sydney will hit a local edge node. On first access (cache miss), the edge fetches from \texttt{us-east-1} and caches the image. All subsequent accesses from that edge serve locally with sub-10ms latency instead of 200--300ms round-trip to Virginia. Tradeoffs: cache invalidation becomes necessary when images change (set appropriate Cache-Control headers and use versioned keys or CloudFront invalidation APIs); CDN adds egress cost (\$0.09/GB from S3 to CloudFront, then \$0.0085--0.02/GB from CloudFront to users depending on region), but this is typically lower than the S3 egress rate for high-traffic content; cold objects on rarely-visited edge nodes will still experience cache-miss latency until they are cached.

\medskip
\textbf{Answer 4.} Move logs to S3 Glacier Deep Archive at \$0.00099/GB-month. For 7-year retention, calculate storage cost: at \$0.023/GB-month (Standard), \$50K/month implies roughly 2.2 PB stored. At \$0.00099/GB-month (Deep Archive), the same data costs approximately \$2,200/month---a 95\% reduction. The implication your team must accept: retrieval from Deep Archive takes 12--48 hours and costs an additional \$0.0025/GB. For annual compliance audits where a week's notice is typical, this retrieval latency is acceptable. Set a lifecycle policy: keep logs in S3 Standard for 30 days (for operational access), transition to Standard-IA for months 2--6, then to Glacier Deep Archive after 6 months. This provides fast operational access while minimizing long-term cost.

\medskip
\textbf{Answer 5.} As of December 2020, S3 provides strong read-after-write consistency for all GET, PUT, and DELETE operations on individual objects. A GET immediately following a PUT is guaranteed to return the new data. If two clients simultaneously PUT to the same key, S3 accepts both writes as separate requests. The last write wins---whichever PUT request completes last (as serialized by S3's internal metadata layer) becomes the current object. The earlier write's data is overwritten and lost. This is a problem in designs where concurrent writes to the same key can occur. Solutions: (1) avoid concurrent writes by design (give each upload a unique key, e.g., including a UUID or timestamp); (2) use DynamoDB conditional writes to coordinate writers before they upload; (3) use S3 versioning, so both writes create distinct versions and application logic can merge or select the correct one. In most well-designed object storage systems, concurrent writes to the same key are avoided architecturally rather than handled at the storage layer.

\end{answerbox}

% ─────────────────────────────────────────────────────────────────────────────
\section{What's Next}
% ─────────────────────────────────────────────────────────────────────────────

\begin{insightbox}
\textbf{Next lesson: L1-10 --- Distributed Systems Fundamentals.}

You have now seen the storage layer in depth. A consistent theme across this lesson---and across every lesson in Layer 1---is the challenge of maintaining consistency across multiple machines. S3 grapples with it (the 2020 consistency upgrade was a multi-year engineering effort). Databases grapple with it. Every horizontally scaled system does. L1-10 confronts this directly: CAP theorem, consistency models (strong, eventual, causal), consensus algorithms (Raft, Paxos), and the concrete failure modes that distributed systems experience in production.

Understanding distributed systems fundamentals is what separates a design that naively assumes ``just add more servers'' from one that reasons carefully about what can go wrong when machines fail, networks partition, and clocks drift. L1-10 provides the vocabulary and the mental models you need to have that reasoning in an interview without handwaving. After L1-10, you will have the complete Layer 1 foundation: networking, push delivery, DNS/CDN, load balancing, relational and NoSQL databases, caching, message queues, object storage, and distributed fundamentals. That unlocks Gate 2 and the first mock interview session.
\end{insightbox}

\end{document}
